\chapter{Bad Taste in the Mouth (Dysgeusia)}

\section*{Definition}
A persistent abnormal or distorted sense of taste, ranging from a metallic, salty, bitter, or sour taste in the mouth without an external stimulus, caused by disorders of the taste buds, olfactory system, or systemic conditions.

\section*{When to See a Doctor}

\begin{itemize}
 \item Persistent bad taste with weight loss, difficulty swallowing, or any neurological symptoms (facial droop, weakness, speech changes)
\end{itemize}

\section*{How It Develops}
Taste is mediated by taste buds on the tongue, palate, pharynx, and epiglottis, innervated by cranial nerves VII (chorda tympani), IX (glossopharyngeal), and X (vagus). Dysgeusia results from damage to taste receptors, alteration of saliva composition, or disruption of central taste pathways. Olfactory dysfunction (as in sinusitis or COVID-19) significantly alters flavor perception. Systemic conditions change saliva composition affecting taste receptor microenvironment.

\section*{Causes}
\begin{itemize}
 \item Oral infections (gingivitis, periodontitis, oral thrush, sinusitis)
 \item Medications (metronidazole, clarithromycin, ACE inhibitors, metformin, lithium, chemotherapy)
 \item Dental problems (poor oral hygiene, dental caries, abscesses, amalgam fillings)
 \item GERD (regurgitation of acidic stomach contents)
 \item Sjogren syndrome or xerostomia (reduced saliva alters taste)
 \item Neurological disorders (Bell palsy, multiple sclerosis, stroke affecting taste pathways)
 \item COVID-19 infection (taste and smell dysfunction)
 \item Nutritional deficiencies (zinc, vitamin B12, copper)
 \item Smoking and tobacco use
 \item Radiation therapy to head and neck (damage to taste buds and salivary glands)
\end{itemize}

\section*{Related Symptoms}
\begin{itemize}
 \item Bad breath
 \item Dry mouth
 \item Loss of taste
 \item Oral thrush
 \item Sinus pressure
\end{itemize}

\section*{Medical Evaluation}
\begin{itemize}
 \item Your doctor will examine the oral cavity for signs of gingivitis, periodontitis, oral thrush, or dental abscesses and review your current medications.
 \item Laboratory testing includes complete blood count, iron studies, zinc and vitamin B12 levels, and blood glucose or HbA1c to screen for nutritional deficiencies and diabetes.
 \item MRI or CT of the brain with contrast is indicated if a neurological cause (stroke, multiple sclerosis, cranial nerve lesion) is suspected based on associated symptoms.
 \item Referral to an otolaryngologist or neurologist is arranged for persistent dysgeusia after initial workup, especially if olfactory dysfunction is present.
\end{itemize}

\section*{Treatment Options}
\begin{itemize}
 \item Treatment of underlying infections with appropriate antibiotics, antifungals, or antiviral medications resolves taste disturbance when infection is the cause.
 \item Switching or adjusting medications suspected of causing dysgeusia (e.g., metronidazole, ACE inhibitors) may be trialled in consultation with the prescribing physician.
 \item Zinc supplementation (50 mg daily) or alpha-lipoic acid may be prescribed for idiopathic dysgeusia, though evidence of benefit is modest.
 \item Management of GERD with proton pump inhibitors, treatment of Sjogren-related dry mouth, and correction of nutritional deficiencies address secondary causes.
\end{itemize}

\section*{Self-Care and Home Management}
\begin{itemize}
 \item Maintain good oral hygiene: brush twice daily, floss, and clean the tongue dorsum to reduce bacterial load in the mouth.
 \item Rinse the mouth with a mild saltwater solution (1/2 teaspoon salt in a glass of warm water) several times daily to freshen the mouth.
 \item Drink plenty of water and chew sugar-free gum to stimulate saliva flow if dry mouth is contributing to the bad taste.
 \item Keep a diary noting when the bad taste occurs, what foods or medications precede it, and any associated symptoms to share with your doctor.
\end{itemize}

\section*{References}
\begin{thebibliography}{99}
\bibitem{bad_taste_in_the_mouth:ref1} Henkin RI. ``Taste dysfunction.'' \textit{N Engl J Med}. 2024;390(6):543-553.
\bibitem{bad_taste_in_the_mouth:ref2} Malaty J, Malaty IAC. ``Smell and taste disorders in primary care.'' \textit{Am Fam Physician}. 2013;88(12):852-859.
\bibitem{bad_taste_in_the_mouth:ref3} Doty RL, Bromley SM. ``Disorders of smell and taste.'' In: \textit{Cummings Otolaryngology}, 7th ed. Elsevier, 2021.
\bibitem{bad_taste_in_the_mouth:ref4} Mann NM. ``Management of smell and taste disorders.'' \textit{Continuum (Minneap Minn)}. 2018;24(5):1408-1425.
\bibitem{bad_taste_in_the_mouth:ref5} UpToDate. ``Taste and olfactory disorders in adults.'' Wolters Kluwer, 2023.
\end{thebibliography}
